\paragraph{Effect of vanishing boundary density on support recovery.}
For $\beta\ge0$, define
\[
p_\beta(x,y)=\frac{(1-r(x,y))^\beta}{Z_\beta}
\mathbf 1_{\{r(x,y)<1\}},\qquad r(x,y):=\sqrt{(x-1)^2+y^2},
\]
where
\[
Z_\beta=2\pi\int_0^1r(1-r)^\beta\,dr
=\frac{2\pi}{(\beta+1)(\beta+2)}.
\]
We consider $\beta\in\{0,2,4\}$. The case $\beta=0$ is uniform on the unit
disk, whereas $\beta=2$ and $\beta=4$ vanish increasingly rapidly at its
boundary. All three distributions have support equal to the closed unit disk.

We fix the support, the dynamics $\dot x=2x,\ \dot y=0$, and $T=0.5$ in every
condition, so the exact reachable set is always the ellipse
\[
S_T=\{(X,Y):(X-e^{2T})^2/e^{4T}+Y^2\le1\}.
\]
Figure~\ref{fig:density-boundary-ablation} compares convex-hull,
packing-ball-union, and Christoffel-function estimators over identical sample
budgets and random-seed indices. Larger $\beta$ makes boundary observations
rarer and increases finite-sample Hausdorff error, despite the unchanged
support. Thus, full support alone does not provide a uniform finite-sample
Hausdorff recovery rate. The packing and Christoffel results also display
estimator-dependent outer deviation.

\begin{figure}[htbp]
    \centering
    \includegraphics[width=\linewidth]{CoRL_2026/fig/density_boundary.png}
    \caption{Effect of vanishing boundary density on reachable-set
    approximation. The initial support is the unit disk and the dynamics,
    time horizon, and exact terminal ellipse are identical in every
    experiment. Samples are drawn from
    $p_\beta(x,y)\propto(1-\sqrt{(x-1)^2+y^2})^\beta$ with
    $\beta\in\{0,2,4\}$. The case $\beta=0$ is uniform, while larger values
    of $\beta$ vanish more rapidly near the boundary. Although all
    distributions have the same support, stronger boundary-density decay
    leads to larger finite-sample Hausdorff error. Curves show means and
    shaded regions show 5th--95th percentiles over 50 seeds.}
    \label{fig:density-boundary-ablation}
\end{figure}

\paragraph{Implementation details.}
We sample $R_\beta\sim\mathrm{Beta}(2,\beta+1)$ and
$\Theta\sim\mathrm{Unif}[0,2\pi)$, then set
$(X,Y)=(R_\beta\cos\Theta,R_\beta\sin\Theta)$. The packing estimator uses the
fixed radius $h=0.05$ for every sample size and density. The Christoffel estimator uses polynomial degree
six, regularization $10^{-6}$, the same empirical maximum threshold rule, and
the same evaluation grid for all three densities.
